% ============================================================
% f2012_jiangsu_q04_flowchart.tex
% 2012 江苏卷第 4 题 算法流程图（TikZ 矢量重绘）
% 开始 → k←1 → 判断 k^2-5k+4>0；否：k←k+1 回代；是：输出 k → 结束
% ============================================================
\begin{tikzpicture}[
  x=1cm, y=1cm, line width=0.6pt, font=\footnotesize, >=Stealth,
  every node/.style={inner sep=3pt},
  proc/.style={rectangle, draw, minimum width=2.4cm, minimum height=0.62cm},
  io/.style={trapezium, draw, trapezium left angle=70, trapezium right angle=110,
             minimum width=2.8cm, minimum height=0.66cm},
  dec/.style={diamond, draw, aspect=2.6, minimum width=3.4cm, minimum height=0.95cm},
  term/.style={rectangle, rounded corners=8pt, draw, minimum width=2.0cm,
               minimum height=0.58cm}
]

  \node[term] (start) at (0,0) {开始};
  \node[proc] (init)  at (0,-1.35) {$k\leftarrow1$};
  \node[dec]  (judge) at (0,-2.95) {$k^2-5k+4>0$};
  \node[proc] (inc)   at (4.10,-2.95) {$k\leftarrow k+1$};
  \node[io]   (out)   at (0,-4.45) {输出 $k$};
  \node[term] (end)   at (0,-5.75) {结束};

  \draw[->] (start) -- (init);
  \draw[->] (init) -- (judge);
  \draw[->] (judge) -- (out);
  \node at (0.55,-3.62) {Y};
  \draw[->] (out) -- (end);
  \draw[->] (judge.east) -- node[above] {N} (inc.west);

  % 否：k←k+1 之后回到判断框之前
  \draw[->] (inc.north) -- (4.10,-1.95) -- (0,-1.95);

\end{tikzpicture}
